\documentclass[11pt]{article}
\usepackage[utf8]{inputenc}
\usepackage[margin=1in]{geometry}
\usepackage{amsmath}
\usepackage{amssymb}
\usepackage{xcolor}
\usepackage{listings}
\usepackage{tcolorbox}
\usepackage{hyperref}

% Define colors for code
\definecolor{codegreen}{rgb}{0,0.6,0}
\definecolor{codegray}{rgb}{0.5,0.5,0.5}
\definecolor{codepurple}{rgb}{0.58,0,0.82}
\definecolor{backcolour}{rgb}{0.95,0.95,0.92}

\lstset{
    backgroundcolor=\color{backcolour},   
    commentstyle=\color{codegreen},
    keywordstyle=\color{magenta},
    numberstyle=\tiny\color{codegray},
    stringstyle=\color{codepurple},
    basicstyle=\ttfamily\footnotesize,
    breaklines=true,
    captionpos=b,                    
    keepspaces=true,                 
    numbers=left,                    
    showspaces=false,                
    showstringspaces=false,
    showtabs=false,                  
    tabsize=2
}

\title{\textbf{Backpropagation from Scratch: A Beginner's Guide}}
\author{A Deep Dive into the Micrograd Engine}
\date{}

\begin{document}

\maketitle

\section{The Core Concept: The ``Paper Trail''}
In standard programming, if you calculate $c = a + b$, the computer gives you the result and immediately forgets how it got there. In Deep Learning, we need the computer to remember.

Think of our \texttt{Value} object as a piece of paper that records:
\begin{enumerate}
    \item \textbf{The Data:} The current number (e.g., 5.0).
    \item \textbf{The Ancestry:} Which other pieces of paper were used to calculate this one?
    \item \textbf{The Operation:} Was it addition? Multiplication? 
    \item \textbf{The Gradient:} How much does this specific number affect the final result?
\end{enumerate}

\subsection{The Chain Rule (The Secret Sauce)}
The entire engine relies on the \textbf{Chain Rule}. If $z$ depends on $y$, and $y$ depends on $x$, then the influence of $x$ on $z$ is:
$$\frac{dz}{dx} = \frac{dz}{dy} \cdot \frac{dy}{dx}$$
In our code, \texttt{self.grad} represents $\frac{dL}{d(\text{self})}$, where $L$ is the final output.

\section{The Implementation: Line-by-Line}

\subsection{1. The Constructor}
We initialize every number with a gradient of $0$. We don't know the influence yet!
\begin{lstlisting}[language=Python]
class Value:
    def __init__(self, data, _children=(), _op='', _label=''):
        self.data = data
        self.grad = 0.0 # Initially, this number has no known impact
        self._prev = set(_children) # The "parents"
        self._backward = lambda: None # Empty "recipe" for math
\end{lstlisting}

\subsection{2. Addition: Passing the Influence}
When we add two numbers, $out = a + b$, the derivative of $out$ with respect to $a$ is just $1$. So, $a$ simply inherits whatever gradient $out$ has.
\begin{lstlisting}[language=Python]
def __add__(self, other):
    other = other if isinstance(other, Value) else Value(other)
    out = Value(self.data + other.data, (self, other), '+')
    
    def _backward():
        self.grad += 1.0 * out.grad
        other.grad += 1.0 * out.grad
    out._backward = _backward
    return out
\end{lstlisting}
\textit{Note: We use \texttt{+=} because if a variable is used twice, its gradients should accumulate (add up).}

\subsection{3. Multiplication: The Switcheroo}
If $out = a \cdot b$, the derivative with respect to $a$ is $b$.
\begin{lstlisting}[language=Python]
def __mul__(self, other):
    other = other if isinstance(other, Value) else Value(other)
    out = Value(self.data * other.data, (self, other), '*')
    
    def _backward():
        self.grad += other.data * out.grad
        other.grad += self.data * out.grad
    out._backward = _backward
    return out
\end{lstlisting}

\section{The ``Brain'' of the Operation: Topological Sort}
We cannot calculate a parent's gradient until we have calculated the child's gradient. To do this, we lay all our ``pieces of paper'' out in a line such that every parent comes \textbf{after} its children. Then, we walk through that line backwards.

\section{Full Code Listing}
\begin{lstlisting}[language=Python]
import math

class Value:
    def __init__(self, data, _children=(), _op='', label=''):
        self.data = data
        self.grad = 0.0
        self._backward = lambda: None
        self._prev = set(_children)
        self._op = _op
        self.label = label

    def __add__(self, other):
        other = other if isinstance(other, Value) else Value(other)
        out = Value(self.data + other.data, (self, other), '+')
        def _backward():
            self.grad += 1.0 * out.grad
            other.grad += 1.0 * out.grad
        out._backward = _backward
        return out

    def __mul__(self, other):
        other = other if isinstance(other, Value) else Value(other)
        out = Value(self.data * other.data, (self, other), '*')
        def _backward():
            self.grad += other.data * out.grad
            other.grad += self.data * out.grad
        out._backward = _backward
        return out

    def tanh(self):
        x = self.data
        t = (math.exp(2*x) - 1)/(math.exp(2*x) + 1)
        out = Value(t, (self, ), 'tanh')
        def _backward():
            self.grad += (1 - t**2) * out.grad
        out._backward = _backward
        return out

    def backward(self):
        topo = []
        visited = set()
        def build_topo(v):
            if v not in visited:
                visited.add(v)
                for child in v._prev:
                    build_topo(child)
                topo.append(v)
        build_topo(self)
        self.grad = 1.0
        for node in reversed(topo):
            node._backward()
\end{lstlisting}

\section{Step-by-Step Explanation of the Micrograd Notebook}

\subsection{1. Importing Libraries}
\begin{lstlisting}[language=Python]
import numpy as np
import math
import graphviz
import matplotlib.pyplot as plt
\end{lstlisting}
\textbf{Explanation:}
\begin{itemize}
    \item \texttt{numpy (np)}: Helps us work with numbers and arrays easily.
    \item \texttt{math}: Lets us use mathematical functions like exponentials.
    \item \texttt{graphviz}: Used to draw diagrams showing how numbers are connected.
    \item \texttt{matplotlib.pyplot (plt)}: Used to make plots and graphs.
\end{itemize}

\subsection{2. The Value Class: The Heart of Everything}
\begin{lstlisting}[language=Python]
class Value:
    def __init__(self, data, _children=(), _op='', _label=''):
        self.data = data
        self._prev = set(_children)
        self._op = _op
        self._label = _label
        self.grad = 0.0
        self._backward = lambda:None
\end{lstlisting}
\textbf{Explanation:}
\begin{itemize}
    \item \texttt{data}: The actual number we are working with.
    \item \texttt{\_prev}: The numbers that were used to make this one (its parents).
    \item \texttt{\_op}: What operation made this number? (e.g., +, *)
    \item \texttt{\_label}: A name for this number, just for us to keep track.
    \item \texttt{grad}: How much this number affects the final result (used for learning).
    \item \texttt{\_backward}: A recipe for how to send influence (gradients) backward.
\end{itemize}

\subsection{3. Printing and Adding Numbers}
\begin{lstlisting}[language=Python]
    def __repr__(self):
        return f"Value(data={self.data})"

    def __add__(self, other):
        out = Value(self.data + other.data, (self, other), "+")
        def _backward():
            self.grad += 1.0*out.grad
            other.grad += 1.0*out.grad
        out._backward = _backward
        return out
\end{lstlisting}
\textbf{Explanation:}
\begin{itemize}
    \item \texttt{\_\_repr\_\_}: Lets us print the Value nicely.
    \item \texttt{\_\_add\_\_}: Lets us add two Value numbers together. It also sets up how to send influence backward for learning.
\end{itemize}

\subsection{4. Multiplying Numbers}
\begin{lstlisting}[language=Python]
    def __mul__(self, other):
        out = Value(self.data * other.data, (self, other), '*')
        def _backward():
            self.grad += other.data*out.grad
            other.grad += self.data*out.grad
        out._backward = _backward
        return out
\end{lstlisting}
\textbf{Explanation:}
\begin{itemize}
    \item \texttt{\_\_mul\_\_}: Multiplies two numbers and sets up the backward recipe for learning.
\end{itemize}

\subsection{5. Tanh Activation: Squashing the Number}
\begin{lstlisting}[language=Python]
    def tanh(self):
        n = self.data
        t = ((math.exp(2*n)-1)/(math.exp(2*n)+1))
        out=Value(t,(self, ),"tanh")
        def _backward():
            self.grad += (1-t**2)*out.grad
        out._backward = _backward
        return out
\end{lstlisting}
\textbf{Explanation:}
\begin{itemize}
    \item \texttt{tanh}: Squashes the number between -1 and 1. This is useful for neural networks to keep numbers manageable.
\end{itemize}

% The backward recipe uses the derivative of tanh:
\begin{equation*}
1 - \tanh^2(x)
\end{equation*}

\subsection{6. Backward Pass: Sending Influence Backwards}
\begin{lstlisting}[language=Python]
    def backward(self):
        topo = []
        visited = set()
        def build_topo(v):
            if v not in visited:
                visited.add(v)
                for child in v._prev:
                    build_topo(child)
                topo.append(v)
        build_topo(self)
        self.grad = 1.0
        for node in reversed(topo):
            node._backward()
\end{lstlisting}
\textbf{Explanation:}
\begin{itemize}
    \item This function figures out the order to send gradients backward so every number gets its correct influence.
    \item It starts from the final result and works backward through all the numbers.
\end{itemize}

\subsection{7. Creating Input Values}
\begin{lstlisting}[language=Python]
a = Value(2.0, _label='a')
b = Value(-3.0, _label='b')
c = Value(10.0, _label='c')
f = Value(-2.0, _label='f')
\end{lstlisting}
\textbf{Explanation:}
\begin{itemize}
    \item We create four numbers: a, b, c, and f. Each is wrapped in a Value object so we can track its history and influence.
    \item \verb|_label| is just a name to help us keep track of which number is which.
\end{itemize}

\subsection{8. Forward Pass: Calculating the Output}
\begin{lstlisting}[language=Python]
e = a * b
e._label = "e"
d = e + c
d._label = 'd'
L = d * f
L._label = 'L'
\end{lstlisting}
\textbf{Explanation:}
\begin{itemize}
    \item We combine our numbers using multiplication and addition, just like you would in math class.
    \item Each step creates a new Value object, which remembers how it was made.
    \item \texttt{L} is our final output, the result of all these operations.
\end{itemize}

\subsection{9. Visualizing the Computation Graph}
\begin{lstlisting}[language=Python]
from graphviz import Digraph

def trace(root):
    nodes, edges = set(), set()
    def build(v):
        if v not in nodes:
            nodes.add(v)
            for child in v._prev:
                edges.add((child, v))
                build(child)
    build(root)
    return nodes, edges

def draw_dot(root):
    dot = Digraph(format='svg', graph_attr={"rankdir": "LR"})
    nodes, edges = trace(root)
    for n in nodes:
        uid = str(id(n))
        dot.node(name=uid, label="{%s | data %.4f | grad %.4f}" % (n._label, n.data, n.grad), shape='record')
        if n._op:
            dot.node(name=uid + n._op, label=n._op)
            dot.edge(uid + n._op, uid)
    for n1, n2 in edges:
        dot.edge(str(id(n1)), str(id(n2)) + n2._op)
    return dot
\end{lstlisting}
\textbf{Explanation:}
\begin{itemize}
    \item This code draws a picture showing how all our numbers are connected.
    \item Each box shows a number, its value, and its gradient.
    \item Arrows show which numbers were used to make other numbers.
\end{itemize}

\subsection{10. Drawing the Graph}
\begin{lstlisting}[language=Python]
draw_dot(L)
\end{lstlisting}
\textbf{Explanation:} This draws the computation graph for our final output L, so we can see the "paper trail" of calculations.

\subsection{11. Gradient Descent: Learning by Nudging}
\begin{lstlisting}[language=Python]
learning_rate = 0.01

a.data += learning_rate * a.grad
b.data += learning_rate * b.grad
c.data += learning_rate * c.grad
f.data += learning_rate * f.grad

# Recompute forward pass
e = a * b
d = e + c
L = d * f

print(L.data)
\end{lstlisting}
\textbf{Explanation:}
\begin{itemize}
    \item \texttt{learning\_rate}: A small number that controls how much we change our values to improve the result.
    \item We "nudge" each number in the direction that makes the output better (lower or higher, depending on our goal).
    \item After nudging, we recalculate everything to see the new result.
    \item \texttt{print(L.data)} shows us the new output after learning.
\end{itemize}

\subsection{12. Manual Gradients: Setting Influence by Hand}
\begin{lstlisting}[language=Python]
L.grad = 1.0
d.grad = -2.0
f.grad = 4.0
e.grad = -2.0
c.grad = -2.0
a.grad = 6.0
b.grad = -4.0
\end{lstlisting}
\textbf{Explanation:}
\begin{itemize}
    \item Here, we set the gradients (influence) for each number by hand, as if we calculated them ourselves.
    \item This is useful for understanding how each number affects the final result.
\end{itemize}

\subsection{13. Plotting the Tanh Function}
\begin{lstlisting}[language=Python]
plt.plot(np.arange(-5,5,0.2),np.tanh(np.arange(-5,5,0.2)))
plt.grid()
\end{lstlisting}
\textbf{Explanation:}
\begin{itemize}
    \item This draws a graph of the tanh function, showing how it squashes numbers between -1 and 1.
    \item The grid helps us see the shape of the curve.
\end{itemize}

\subsection{14. Building a Simple Neuron}
\begin{lstlisting}[language=Python]
#inputs X1,X2
x1 = Value(2.0, _label="x1")
x2 = Value(0.0, _label="x2")
#weights 
w1 = Value(-3.0, _label="w1")
w2 = Value(1.0, _label="w2")
#bias of the neuron
b = Value(6.8813735, _label="bias")
x1w1 = x1 * w1; x1w1._label = "x1*w1"
x2w2 = x2 * w2; x2w2._label = "x2*w2"
x1w1x2w2 = x1w1 + x2w2; x1w1x2w2._label = "x1*w1 + x2*w2"
n = x1w1x2w2 + b; n._label = "n"
o = n.tanh(); o._label = "o"
\end{lstlisting}
\textbf{Explanation:}
\begin{itemize}
    \item \texttt{x1}, \texttt{x2}: These are the inputs to our neuron (like features in a dataset).
    \item \texttt{w1}, \texttt{w2}: These are weights, which tell us how important each input is.
    \item \texttt{b}: The bias, which shifts the output up or down.
    \item We multiply each input by its weight, add them together, and add the bias.
    \item \texttt{n}: The total input to the neuron before activation.
    \item \texttt{o}: The output after applying the tanh activation, squashing the result between -1 and 1.
\end{itemize}

\subsection{15. Visualizing the Neuron}
\begin{lstlisting}[language=Python]
draw_dot(o)
\end{lstlisting}
\textbf{Explanation:} This draws the computation graph for the neuron's output, showing all the steps and connections.

\end{document}